%-*-latex-*-
%-*-latex-*-
\newcommand\COURSE{ciss350}
\newcommand\ASSESSMENT{q10401}
\newcommand\ASSESSMENTTYPE{Quiz}
\newcommand\POINTS{\textwhite{xxx/xxx}}

\input{myquizpreamble}
\input{yliow}
\input{\COURSE}
%-*-latex-*-

%-*-latex-*-

%-*-latex-*-
\renewcommand\TITLE{\ASSESSMENTTYPE \ \ASSESSMENT}

\renewcommand\EMAIL{}
\renewcommand\AUTHOR{eagraham1@cougars.ccis.edu}

\textwidth=6in
\begin{document}
\topmatterthree

%------------------------------------------------------------------------------
\nextq
You want to sort an array of size $n$ (which is $\geq 1000000000$)
of integer value in the range 0..9.
You want the fastest sorting algorithm.
Complete the following function to essentially perform a counting sort on \verb!x!.
\\
\ANSWER
\begin{answercode}
// Sort x[0..n-1]
void sort(int x[], int n)
{
  int count[10] = {0};

  for(int i = 0; i < n; ++i)
  {
    count[x[i]]++;
  }

  int index = 0;
  for(int i = 0; i < 10; i++)
  {
    while(count[i] > 0)
    {
      x[index++];
      count[i]--;
    }
  }
}
\end{answercode}

%------------------------------------------------------------------------------
\nextq
You want to sort an array of size $n$ (which is $\geq 1000000000$).
of integer value in the range 100000..100020.
The device is memory-constrained -- besides the input array,
you can afford to use at most
the equivalent of 100 integers for your computations (for variables).
Complete the following function to essentially perform a counting sort on \verb!x!.
\\
\ANSWER
\begin{answercode}
// Sort x[0..n-1]
void sort(int x[], int n)
{
  const int min = 100000;
  const int range = 21;

  count[range] = 0;

  for (int i = 0; i < n; ++i)
  {
    count[x[i] - MIN]++;
  }

  int index = 0;
  for (int i = 0; i < RANGE; ++i)
  {
    while (count[i] > 0)
    {
      x[index++] = i + MIN;
      count[i]--;
    }
  }
}
\end{answercode}

%------------------------------------------------------------------------------
\nextq
You want to sort an array of size $n$ (which is $\geq 1000000000$)
and contains
double values -5.0, -4.5, -4.0, -3.5, -3.0, ..., -0.5, 0.0, 0.5, ..., 4.5, 5.0.
The device is memory-constrained -- besides the input array,
you can afford to use at most
the equivalent of 100 integers for your computations (for variables).
Complete the following function to essentially perform a counting sort on \verb!x!.
\\
\ANSWER
\begin{answercode}
// Sort x[0..n-1]
void sort(double x[], int n) // correction: int x[] should be double x[].
{
  const int RANGE = 21; // -5.0 to 5.0 step 0.5
  int count[RANGE] = {0};

  for (int i = 0; i < n; ++i)
  {
    int idx = (int)((x[i] + 5.0) / 0.5);
    count[idx]++;
  }

  int index = 0;
  for (int i = 0; i < RANGE; ++i)
  {
    double value = -5.0 + i * 0.5;
    while (count[i] > 0)
    {
      x[index++] = value;
      count[i]--;
    }
  }
}
\end{answercode}

%------------------------------------------------------------------------------
\nextq
Given a string, you want to return the character which appears most
frequently in the string.
For instance if the most frequently occurring character in
\verb!"hello world"! is \verb!'l'!.
Use the idea from \verb!countingsort! to count the number of times
a character occurs in a string \verb!s! and then perform a max
computation to compute the most commonly occurring character in the given
string.
\\
\ANSWER
\begin{answercode}
char most_common_char(char s[])
{
  int count[256]; // count[i] is the number of times char(i) occurs in s

  for (int i = 0; s[i] != '\0'; ++i)
  {
    count[(int)s[i]]++;
  }

  int max = 0;
  char result = s[0];

  for (int i = 0; i < 256; ++i)
  {
    if (count[i] > max)
    {
      max = count[i];
      result = (char)i;
    }
  }

  return result;
}
\end{answercode}

%------------------------------------------------------------------------------
%-*-latex-*-
\newpage
\input{instructions}
\end{document}

